\documentclass[12pt]{article}
\usepackage[margin=0.75in]{geometry}
\usepackage{lmodern}
\usepackage{fancyhdr}
\usepackage{titlesec}
\usepackage{enumitem}
\usepackage{xcolor}
\usepackage[hidelinks]{hyperref}

\setlength{\parindent}{1.1em}
\setlength{\parskip}{0.35em}
\setlength{\headheight}{15pt}
\titleformat{\section}{\large\bfseries\scshape}{}{0pt}{}
\titleformat{\subsection}{\normalsize\bfseries}{}{0pt}{}
\pagestyle{fancy}
\fancyhf{}
\fancyfoot[C]{\thepage}
\renewcommand{\headrulewidth}{0.4pt}

\fancyhead[L]{Whately Elements of Logic}
\fancyhead[R]{Teacher Guide}
\begin{document}
\begin{center}
{\LARGE\bfseries Teacher Guide: Week 3}\par\vspace{0.25em}
{\large\scshape What We Say and What We Mean: Making Clear Claims}\par\vspace{0.4em}
{10th Grade CUE Unit Based on Richard Whately's \textit{Elements of Logic}}\par\vspace{0.6em}
\rule{0.82\textwidth}{0.4pt}
\end{center}

\section*{Week 3 Overview}
Unit 3 moves from terms (Unit 2) to propositions—the statements we make when we combine terms. Students learn to identify the subject, predicate, quality, and quantity of claims, and to translate ordinary language into the four standard forms. The unit emphasizes precision in stating what is being claimed about how much of a subject and in what way. This foundation is essential for the study of syllogisms in later weeks.

\section*{Essential Question}
How can we state claims clearly enough that others know exactly what we are asserting?

\section*{Student-Facing Materials}
\begin{itemize}[leftmargin=*]
\item \textit{What We Say and What We Mean: Making Clear Claims} — student reading handout.
\item \textit{Week 3 Argument Lab: Identifying and Clarifying Propositions} — practice in identifying quality, quantity, distribution, and translating ordinary language.
\end{itemize}

\section*{Learning Goals}
Students will be able to:
\begin{enumerate}[leftmargin=*]
\item Define a proposition and distinguish it from questions, commands, and exclamations.
\item Identify the subject, predicate, and copula of a proposition.
\item Distinguish affirmative from negative propositions (quality).
\item Distinguish universal from particular propositions (quantity).
\item Recognize the four standard forms (A, E, I, O).
\item Determine whether terms are distributed or undistributed.
\item Translate common English expressions into standard logical form.
\end{enumerate}

\section*{Key Vocabulary}
proposition, subject, predicate, copula, quality, quantity, universal, particular, distribution, affirmative, negative

\section*{Weekly Lesson Sequence}

\subsection*{Day 1: What Is a Proposition?}
\begin{itemize}[leftmargin=*]
\item Begin with examples of statements that can be true or false versus questions and commands.
\item Introduce the idea that logic is concerned with claims that can be evaluated.
\item Read the first half of the student reading.
\item Students identify five propositions from everyday language and explain why they qualify.
\end{itemize}

\subsection*{Day 2: Quality and Quantity}
\begin{itemize}[leftmargin=*]
\item Focus on the distinction between affirmative/negative and universal/particular.
\item Use multiple examples on the board. Have students classify each one.
\item Complete the first two sections of the Argument Lab in pairs.
\end{itemize}

\subsection*{Day 3: The Four Forms and Distribution}
\begin{itemize}[leftmargin=*]
\item Introduce the A, E, I, O labels.
\item Practice determining distribution of terms.
\item Students work through Part 2 of the Lab and explain their reasoning.
\end{itemize}

\subsection*{Day 4: Translating Ordinary Language}
\begin{itemize}[leftmargin=*]
\item This is often the most difficult skill. Spend significant time on translation.
\item Work through Part 3 as a class, then have students complete Part 4 independently.
\item Discuss why translation matters for clear argument.
\end{itemize}

\subsection*{Day 5: Window Notes and Review}
\begin{itemize}[leftmargin=*]
\item Students complete the window note task from the Lab.
\item Class discussion: How does making quantity and quality explicit change the way an argument can be supported or challenged?
\item Exit task: Write one claim from the unit in two different logical forms and explain the difference.
\end{itemize}

\section*{Notes on Examples}
The student reading deliberately provides numerous examples of each concept. Students benefit from seeing the same idea illustrated in several different contexts. The translation exercises in the Lab are especially important; many students initially struggle to move between ordinary language and standard form. Encourage them to ask “What exactly is being claimed about every member, some members, or no members of the subject?”

\section*{Connection to Future Units}
This unit prepares students for syllogistic reasoning. In the next unit, students will learn to evaluate whether a conclusion follows from two premises. That evaluation depends entirely on the ability to identify the quality, quantity, and distribution of the propositions involved.
\end{document}